\clearpage
\setcounter{page}{1}
\maketitlesupplementary


\section{Theoretical Foundations for BRIDGE}\label{app:theory}

\noindent
To formalize the problem of imbalanced representation learning, we introduce several key definitions and theorems.

\begin{definition}[Image and Dataset Spaces]
Let $\mathcal{X}$ be the space of all possible images, and $D = \{x_1, x_2, \ldots, x_n\} \subset \mathcal{X}$ be our finite dataset.
\end{definition}

\begin{definition}[Concepts]
Let $\mathcal{C} = \{C_1, C_2, \ldots, C_m\}$ be a partition of $\mathcal{X}$ representing distinct semantic concepts, where each $C_i \subset \mathcal{X}$. We denote the restriction of dataset $D$ to concept $C_i$ as $D_i = D \cap C_i$.
\end{definition}

\begin{definition}[Information Density]
The information density $\mathcal{I}(C_i)$ of a concept $C_i$ is defined as the minimum number of examples needed to represent it effectively. Formally:
\begin{equation}
\mathcal{I}(C_i) = \min\{|S| : S \subset C_i \text{ and } \mathcal{L}(f_S, C_i) \leq \epsilon\}
\end{equation}
where $f_S$ is a representation learned from subset $S$, and $\mathcal{L}(f_S, C_i)$ is a loss function measuring how well $f_S$ represents concept $C_i$.
\end{definition}

\begin{definition}[Balanced Dataset]
A dataset $D$ is $\delta$-balanced with respect to concepts $\mathcal{C}$ if for all concepts $C_i \in \mathcal{C}$:
\begin{equation}
\frac{|D_i|}{\mathcal{I}(C_i)} \geq \delta
\end{equation}
where $\delta \in (0,1]$ is a balance factor. When $\delta = 1$, the dataset has sufficient examples for each concept.
\end{definition}

\begin{assumption}[Ideal Embedding]
\label{assump:ideal}
In our initial analysis, we assume that the encoder $f_\theta$ perfectly clusters images by concepts. Formally, there exists a distance threshold $\tau$ such that:
\begin{align}
\forall x_i \in C_i, x_j \in C_j: \|f_\theta(x_i) - f_\theta(x_j)\|_2 < \tau \iff i = j
\end{align}
\end{assumption}

\begin{theorem}[Embedding Density Reflects Concept Representation]
\label{thm:density}
Under Assumption \ref{assump:ideal}, for any two concepts $C_i$ and $C_j$ with $|D_i|/\mathcal{I}(C_i) < |D_j|/\mathcal{I}(C_j)$, the average k-NN distance for points in $D_i$ is greater than for points in $D_j$:
\begin{equation}
\mathbb{E}_{x \in D_i}[d_{kNN}(f_\theta(x), Z)] > \mathbb{E}_{x \in D_j}[d_{kNN}(f_\theta(x), Z)]
\end{equation}
\end{theorem}

\begin{proof}
Under Assumption \ref{assump:ideal}, embeddings form distinct clusters by concept. Let $Z_i = \{f_\theta(x) : x \in D_i\}$ be the set of embeddings for concept $C_i$. For any $x \in D_i$, its k nearest neighbors must come from the same concept, i.e., from $Z_i$, since the distance to points from other concepts exceeds $\tau$ while distances within the same concept are less than $\tau$.

\noindent
For points uniformly distributed in a d-dimensional space of volume $V$, the expected distance to the k-th nearest neighbor of a random point scales as $\mathcal{O}(V^{1/d} \cdot n^{-1/d})$, where $n$ is the number of points. This scaling relationship follows from the fact that the expected distance between uniformly distributed points in a region of volume $V$ is proportional to $V^{1/d}$, and with $n$ points, the average volume per point is $V/n$.

\noindent
The embeddings of concept $C_i$ occupy a region in the embedding space with some effective volume $V_i$. Since $|Z_i| = |D_i|$, the average k-NN distance for points in $D_i$ scales as:
\begin{equation}
\mathbb{E}_{x \in D_i}[d_{kNN}(f_\theta(x), Z_i)] \propto V_i^{1/d} \cdot |D_i|^{-1/d}
\end{equation}

\noindent
Assuming similar volumes for the concept regions (i.e., $V_i \approx V_j$), we have:
\begin{equation}
\frac{\mathbb{E}_{x \in D_i}[d_{kNN}(f_\theta(x), Z_i)]}{\mathbb{E}_{x \in D_j}[d_{kNN}(f_\theta(x), Z_j)]} \propto \left(\frac{|D_j|}{|D_i|}\right)^{1/d}
\end{equation}

\noindent
Given that $|D_i|/\mathcal{I}(C_i) < |D_j|/\mathcal{I}(C_j)$, and assuming comparable information densities, we have $|D_i| < |D_j|$. Therefore:
\begin{equation}
\begin{aligned}
\mathbb{E}_{x \in D_i}[d_{kNN}(f_\theta(x), Z)] &= \mathbb{E}_{x \in D_i}[d_{kNN}(f_\theta(x), Z_i)] \\
&> \mathbb{E}_{x \in D_j}[d_{kNN}(f_\theta(x), Z_j)] \\
&= \mathbb{E}_{x \in D_j}[d_{kNN}(f_\theta(x), Z)]
\end{aligned}
\end{equation}
\end{proof}

\begin{theorem}[Rank-Based Selection Identifies Underrepresented Concepts]
\label{thm:rank}
Let $O = \{x_i \in D : \text{rank}(d_{kNN}(f_\theta(x_i), Z)) \leq N_{OOD}\}$ be the set of $N_{OOD}$ points with the highest k-NN distances. Under Assumption \ref{assump:ideal}, $O$ contains a higher proportion of points from underrepresented concepts than from well-represented concepts:
\begin{equation}
\frac{|O \cap D_i|}{|D_i|} > \frac{|O \cap D_j|}{|D_j|} \text{ for } \frac{|D_i|}{\mathcal{I}(C_i)} < \frac{|D_j|}{\mathcal{I}(C_j)}
\end{equation}
\end{theorem}

\begin{proof}
From Theorem \ref{thm:density}, points from underrepresented concepts have higher expected k-NN distances. Let $F_i(d)$ be the cumulative distribution function (CDF) of k-NN distances for points in concept $C_i$, i.e., $F_i(d) = \Pr(d_{kNN}(f_\theta(x), Z) \leq d | x \in D_i)$.

\noindent
Under Assumption \ref{assump:ideal} and the uniform distribution assumption within concept regions, the CDF $F_i(d)$ is determined by the point density in the concept region. For concepts $C_i$ and $C_j$ with $|D_i|/\mathcal{I}(C_i) < |D_j|/\mathcal{I}(C_j)$, the lower density of points in concept $C_i$ results in $F_i(d) < F_j(d)$ for all distance thresholds $d$.

\noindent
When selecting the top $N_{OOD}$ points by k-NN distance, there exists some distance threshold $d^*$ such that points with $d_{kNN}(f_\theta(x), Z) > d^*$ are selected. The proportion of points selected from concept $C_i$ is:
\begin{equation}
\frac{|O \cap D_i|}{|D_i|} = 1 - F_i(d^*)
\end{equation}

\noindent
Similarly, for concept $C_j$:
\begin{equation}
\frac{|O \cap D_j|}{|D_j|} = 1 - F_j(d^*)
\end{equation}

\noindent
Since $F_i(d^*) < F_j(d^*)$, we have:
\begin{equation}
\frac{|O \cap D_i|}{|D_i|} = 1 - F_i(d^*) > 1 - F_j(d^*) = \frac{|O \cap D_j|}{|D_j|}
\end{equation}

\noindent
Therefore, underrepresented concepts are proportionally more represented in the selected set $O$.
\end{proof}

\begin{assumption}[Ideal Diffusion]
\label{assump:diffusion}
The diffusion model $g: \mathcal{X} \times \mathbb{R}^r \rightarrow \mathcal{X}$ generates samples that belong to the same concept as the input:
\begin{equation}
\forall x \in C_i, \epsilon \sim \mathcal{N}(0, I_r): g(x, \epsilon) \in C_i
\end{equation}
\end{assumption}

\begin{theorem}[Augmentation Improves Balance]
\label{thm:augment}
Let $D' = D \cup \{g(x, \epsilon_{j,x}) : x \in O, j \in \{1,\ldots,N_{Aug}\}\}$ where $O$ is the set of $N_{OOD}$ points with highest k-NN distances in $Z$, and $N_{Aug}$ is the number of augmentations per point. Under Assumptions \ref{assump:ideal} and \ref{assump:diffusion}, if $N_{OOD}$ and $N_{Aug}$ are sufficiently large, then $D'$ has a higher balance factor than $D$.
\end{theorem}

\begin{proof}
\noindent
From Theorem \ref{thm:rank}, points in $O$ come disproportionately from underrepresented concepts. Let $p_i = |O \cap D_i|/|O|$ be the proportion of points in $O$ from concept $C_i$.

\noindent
Under Assumption \ref{assump:diffusion}, for each $x \in O \cap D_i$, all generated samples $g(x, \epsilon)$ also belong to $C_i$. Therefore, the number of new samples added to concept $C_i$ is $N_{Aug} \cdot |O \cap D_i| = N_{Aug} \cdot p_i \cdot N_{OOD}$.

\noindent
Let the balance factor of the original dataset be:
\begin{equation}
\delta = \min_{i} \frac{|D_i|}{\mathcal{I}(C_i)}
\end{equation}

\noindent
Assuming the minimum is achieved for concept $C_k$ (i.e., $C_k$ is the most underrepresented concept), after augmentation, the new balance factor is:
\begin{equation}
\delta' = \min_{i} \frac{|D_i| + N_{Aug} \cdot |O \cap D_i|}{\mathcal{I}(C_i)}
\end{equation}

\noindent
We need to show that $\delta' > \delta$. Consider the ratio of representation increase for concept $C_i$:
\begin{equation}
\frac{|D_i| + N_{Aug} \cdot |O \cap D_i|}{|D_i|} = 1 + N_{Aug} \cdot \frac{|O \cap D_i|}{|D_i|}
\end{equation}

\noindent
From Theorem \ref{thm:rank}, we know that $\frac{|O \cap D_i|}{|D_i|} > \frac{|O \cap D_j|}{|D_j|}$ for $\frac{|D_i|}{\mathcal{I}(C_i)} < \frac{|D_j|}{\mathcal{I}(C_j)}$. This means that $\frac{|O \cap D_k|}{|D_k|} > \frac{|O \cap D_j|}{|D_j|}$ for all $j$ such that $\frac{|D_k|}{\mathcal{I}(C_k)} < \frac{|D_j|}{\mathcal{I}(C_j)}$.

\noindent
Therefore, the proportional increase for the most underrepresented concept $C_k$ is greater than for other better-represented concepts. This ensures that:
\begin{equation}
\frac{|D_k| + N_{Aug} \cdot |O \cap D_k|}{\mathcal{I}(C_k)} > \frac{|D_k|}{\mathcal{I}(C_k)} = \delta
\end{equation}

\noindent
For any other concept $C_i$, either $\frac{|D_i| + N_{Aug} \cdot |O \cap D_i|}{\mathcal{I}(C_i)} > \frac{|D_i|}{\mathcal{I}(C_i)} \geq \delta$ (if it wasn't the minimizer), or it becomes the new minimizer but with a value greater than $\delta$. Thus, $\delta' > \delta$, meaning the dataset becomes more balanced.
\end{proof}

\begin{theorem}[Convergence of Iterative BRADD]
\label{thm:converge}
\noindent
Under Assumptions \ref{assump:ideal} and \ref{assump:diffusion}, the iterative BRADD process converges to a dataset with balance factor $\delta \rightarrow 1$ as the number of iterations approaches infinity.
\end{theorem}

\begin{proof}
\noindent
Let $\delta^{(t)}$ be the balance factor after $t$ iterations. From Theorem \ref{thm:augment}, we know that $\delta^{(t+1)} > \delta^{(t)}$ for all $t$, provided that $N_{OOD}$ and $N_{Aug}$ are sufficiently large.

\noindent
The sequence $\{\delta^{(t)}\}$ is monotonically increasing and bounded above by 1 (since a perfectly balanced dataset has $\delta = 1$). By the monotone convergence theorem, the sequence converges to a limit $\delta^* \leq 1$.

\noindent
Suppose $\delta^* < 1$. This means there exists at least one concept $C_i$ such that $|D_i^{(\infty)}|/\mathcal{I}(C_i) = \delta^* < 1$, where $D_i^{(\infty)}$ is the limit of the concept subset after infinitely many iterations.

\noindent
Let $C_k$ be such a concept with the minimum ratio: $|D_k^{(\infty)}|/\mathcal{I}(C_k) = \delta^*$. By Theorem \ref{thm:rank}, points from $C_k$ would have higher k-NN distances than points from better-represented concepts, and would be prioritized for augmentation in iteration $t+1$.

\noindent
Specifically, let $p_k^{(t+1)} = |O^{(t+1)} \cap D_k^{(t)}|/|O^{(t+1)}|$ be the proportion of selected points from concept $C_k$ in iteration $t+1$. From Theorem \ref{thm:rank}, we know that $p_k^{(t+1)} > |D_k^{(t)}|/|D^{(t)}|$, meaning concept $C_k$ is overrepresented in the selected set compared to its proportion in the dataset.

\noindent
After augmentation, the new size of concept $C_k$ would be:
\begin{equation}
\begin{aligned}
|D_k^{(t+1)}| &=
|D_k^{(t)}| + N_{Aug} \cdot |O^{(t+1)} \cap D_k^{(t)}|\\
&=
|D_k^{(t)}| + N_{Aug} \cdot p_k^{(t+1)} \cdot N_{OOD}
\end{aligned}
\end{equation}

\noindent
This leads to a new balance factor $\delta^{(t+1)} > \delta^{(t)}$, contradicting the convergence to $\delta^* < 1$.

\noindent
Therefore, $\delta^* = 1$, meaning the iterative process converges to a perfectly balanced dataset.
\end{proof}

\noindent
We extend these results to more realistic settings where embeddings are imperfect and diffusion models might generate out-of-concept samples. Due to space constraints, we provide only the key theorem for practical applications:


\begin{assumption}[Realistic Embedding]
\label{assump:realistic}
The encoder $f_\theta$ approximately clusters images by concepts. Formally, there exist distance thresholds $\tau_1 < \tau_2$ such that:
\begin{align}
\forall x_i \in C_i, x_j \in C_j: \\
\mathrm{Pr}(\|f_\theta(x_i) - f_\theta(x_j)\|_2 < \tau_1 | i = j) \geq 1 - \epsilon_1 \\
\mathrm{Pr}(\|f_\theta(x_i) - f_\theta(x_j)\|_2 > \tau_2 | i \neq j) \geq 1 - \epsilon_2
\end{align}
where $\epsilon_1, \epsilon_2 \in (0,1)$ are small error probabilities.
\end{assumption}

\begin{assumption}[Realistic Diffusion]
\label{assump:realistic_diff}
The diffusion model $g: \mathcal{X} \times \mathbb{R}^r \rightarrow \mathcal{X}$ generates samples that belong to the same concept as the input with high probability:
\begin{equation}
\forall x \in C_i, \epsilon \sim \mathcal{N}(0, I_r): \mathrm{Pr}(g(x, \epsilon) \in C_i) \geq 1 - \epsilon_g
\end{equation}
where $\epsilon_g \in (0,1)$ is a small error probability.
\end{assumption}

\begin{assumption}[Embedding Improvement]
\label{assump:improvement}
The quality of the encoder $f_\theta$ improves with more balanced training data. Formally, if dataset $D'$ has a higher balance factor than $D$, then the encoder $f_{\theta'}$ trained on $D'$ has lower error probabilities $\epsilon_1', \epsilon_2'$ compared to the encoder $f_\theta$ trained on $D$:
\begin{equation}
\epsilon_1' < \epsilon_1 \text{ and } \epsilon_2' < \epsilon_2
\end{equation}
\end{assumption}

\begin{theorem}[Probabilistic Balance Improvement]
\label{thm:prob}
Under Assumptions \ref{assump:realistic}, \ref{assump:realistic_diff}, and \ref{assump:improvement}, the BRADD algorithm improves dataset balance with high probability, and this probability increases as the quality of embeddings improves over iterations.
\end{theorem}

\begin{proof}
Under the realistic embedding assumption, embeddings form approximate clusters by concept. Let $Z_i = \{f_\theta(x) : x \in D_i\}$ be the set of embeddings for concept $C_i$.

For any $x \in D_i$, its k nearest neighbors come from the same concept $C_i$ with probability at least $(1 - \epsilon_1)^k$. With probability at most $1 - (1 - \epsilon_1)^k$, at least one of the k nearest neighbors comes from a different concept.

The expected k-NN distance for points in $D_i$ can be written as:
\begin{equation}
\begin{aligned}
&\mathbb{E}_{x \in D_i}[d_{kNN}(f_\theta(x), Z)] =\\ &(1 - \epsilon_1)^k \cdot \mathbb{E}[d_{kNN}(f_\theta(x), Z_i)]\\
&+ (1 - (1 - \epsilon_1)^k) \cdot \mathbb{E}[d_{kNN}(f_\theta(x), Z \setminus Z_i)]
\end{aligned}
\end{equation}

\noindent
For sufficiently small $\epsilon_1$ and k, the first term dominates, and the expected k-NN distance is approximately:
\begin{equation}
\mathbb{E}_{x \in D_i}[d_{kNN}(f_\theta(x), Z)] \approx \mathbb{E}[d_{kNN}(f_\theta(x), Z_i)] \propto |D_i|^{-1/d}
\end{equation}

\noindent
Given that $|D_i|/\mathcal{I}(C_i) < |D_j|/\mathcal{I}(C_j)$ and assuming similar volumes for the concept regions, we can bound the probability:
\begin{align}
\mathrm{Pr}(\mathbb{E}_{x \in D_i}[d_{kNN}(f_\theta(x), Z)] > \mathbb{E}_{x \in D_j}[d_{kNN}(f_\theta(x), Z)]) \geq 1 - \gamma
\end{align}
where $\gamma = 1 - (1 - \epsilon_1)^{2k} \cdot (1 - \exp(-c \cdot (|D_j|^{-1/d} - |D_i|^{-1/d})^2 \cdot \min(|D_i|, |D_j|)))$ for some constant $c > 0$.

\noindent
When selecting the top $N_{OOD}$ points by k-NN distance (set $O$), the probability that underrepresented concepts are proportionally more represented in $O$ is at least $1 - \gamma'$, where $\gamma'$ accounts for finite sample effects in addition to embedding errors.

\noindent
Specifically, for concepts $C_i$ and $C_j$ with $|D_i|/\mathcal{I}(C_i) < |D_j|/\mathcal{I}(C_j)$:
\begin{equation}
\mathrm{Pr}\left(\frac{|O \cap D_i|}{|D_i|} > \frac{|O \cap D_j|}{|D_j|}\right) \geq 1 - \gamma'
\end{equation}

\noindent
Under the realistic diffusion assumption, for each $x \in O \cap D_i$, the generated sample $g(x, \epsilon)$ belongs to $C_i$ with probability at least $1 - \epsilon_g$. Therefore, the expected number of new samples correctly added to concept $C_i$ is at least $(1 - \epsilon_g) \cdot N_{Aug} \cdot |O \cap D_i|$.

\noindent
Let the balance factor of the original dataset be:
\begin{equation}
\delta = \min_{i} \frac{|D_i|}{\mathcal{I}(C_i)}
\end{equation}

\noindent
Assuming the minimum is achieved for concept $C_k$, after augmentation, the new balance factor is:
\begin{equation}
\delta' = \min_{i} \frac{|D_i| + a_i}{\mathcal{I}(C_i)}
\end{equation}
where $a_i$ is the number of correctly augmented samples added to concept $C_i$.

\noindent
The probability that the balance factor improves ($\delta' > \delta$) depends on:
1. The probability that underrepresented concepts are proportionally more represented in $O$ ($1 - \gamma'$)
2. The probability that the diffusion model correctly preserves concept membership ($1 - \epsilon_g$)
3. The probability that these effects are strong enough to increase the minimum ratio

\noindent
For the most underrepresented concept $C_k$, the expected increase in representation is:
\begin{equation}
(1 - \epsilon_g) \cdot N_{Aug} \cdot |O \cap D_k|
\end{equation}

\noindent
This is proportionally larger than for better-represented concepts with high probability ($1 - \gamma'$). Using Hoeffding's inequality to bound the deviation from expected values, we can establish:
\begin{equation}
\mathrm{Pr}(\delta' > \delta) \geq 1 - \gamma'' - 2\exp(-2N_{OOD}(1-\epsilon_g)^2)
\end{equation}
where $\gamma''$ depends on $\gamma'$ and other parameters.

\noindent
By Assumption \ref{assump:improvement}, as the dataset becomes more balanced after each iteration, the embedding quality improves (i.e., $\epsilon_1$ and $\epsilon_2$ decrease), reducing $\gamma''$ in subsequent iterations. This creates a virtuous cycle where the probability of successful balance improvement increases over time.

\noindent
Let $P_t = \mathrm{Pr}(\delta^{(t+1)} > \delta^{(t)})$ be the probability of improvement at iteration $t$. The probability that all iterations lead to improvements is:
\begin{equation}
\mathrm{Pr}(\forall t: \delta^{(t+1)} > \delta^{(t)}) = \prod_{t=0}^{\infty} P_t
\end{equation}

\noindent
This infinite product converges to a positive value if $\sum_{t=0}^{\infty} (1 - P_t) < \infty$, which holds true if the error probabilities decrease fast enough with iterations (e.g., exponentially). Under these conditions, the BRADD algorithm leads to monotonically increasing balance factors with non-zero probability, converging to a well-balanced dataset.
\end{proof}

\noindent
This proof demonstrates that even with imperfect embeddings and diffusion models, our method can effectively improve dataset balance with high probability, and this probability increases as the training progresses and representations improve.
